\begin{center}
    \textbf{BAB V}\\
    \textbf{KESIMPULAN DAN SARAN}
    \addcontentsline{toc}{chapter}{BAB V. KESIMPULAN DAN SARAN}
\end{center}

\textbf{A. Kesimpulan}\\
\addcontentsline{toc}{section}{A. Kesimpulan}   
\hspace*{1cm}Berdasarkan hasil penelitian pengembangan media pembelajaran matematika berbasis simulasi 3D pada materi bangun ruang sisi datar, dapat disimpulkan bahwa:
\begin{enumerate}
    \item Pengembangan media pembelajaran berbasis simulasi 3D menggunakan model ADDIE sebagai kerangka kerja utama, dengan uraian sebagai berikut:
    \begin{enumerate}[leftmargin=*, label=\arabic*)]
        \item Analisis: peneliti melakukan analisis kebutuhan dengan wawancara kepada pendidik matematika SMP Muhammadiyah 9 Yogyakarta, mengidentifikasi kompetensi dasar, hambatan pembelajaran, dan kebutuhan teknis untuk media simulasi 3D.
        \item Perancangan: peneliti merancang arsitektur sistem, tampilan antarmuka, alur navigasi, materi pembelajaran, storyboard simulasi, serta skenario interaksi pengguna agar sesuai dengan tujuan pembelajaran bangun ruang.
        \item Pengembangan: media dikembangkan sebagai aplikasi website dengan integrasi React, Three.js, model 3D hasil pembuatan di Blender, asset grafis, serta fitur evaluasi seperti isian singkat, soal akhir bab, dan kuis interaktif.
        \item Implementasi: media diuji coba pada peserta didik kelas VIII dengan pelaksanaan pembelajaran berbasis simulasi 3D, pengumpulan respons melalui angket, serta pengamatan terhadap kemudahan pemahaman dan minat belajar siswa.
        \item Evaluasi: hasil validasi oleh ahli materi dan ahli media, serta analisis respons peserta didik, digunakan untuk menilai kelayakan media dan melakukan perbaikan berdasarkan masukan yang diperoleh.
    \end{enumerate}
    \item Kelayak media pembelajaran berbasis simulasi 3D diperoleh dari hasil validasi dan uji coba sebagai berikut:
    \begin{enumerate}[leftmargin=*, label=\arabic*)]
        \item Validasi ahli materi memperoleh skor rata-rata \(4{,}63\) dengan kategori \textbf{Sangat Baik}. Aspek yang dinilai meliputi isi, bahasa, penyajian, dan simulasi 3D.
        \item Validasi ahli media memperoleh skor rata-rata \(4{,}66\) dengan kategori \textbf{Sangat Baik}. Aspek yang dinilai meliputi kualitas kegrafikan, kelayakan bahasa, dan kualitas perangkat lunak.
        \item Uji coba kelas kecil memperoleh skor rata-rata \(4{,}31\) dengan kategori \textbf{Sangat Baik}. Hasil ini menunjukkan bahwa media dinilai menarik dan memudahkan pemahaman oleh peserta didik.
        \item Uji coba kelas besar memperoleh skor rata-rata \(4{,}65\) dengan kategori \textbf{Sangat Baik}. Respons peserta didik menunjukkan bahwa media membantu pemahaman materi dan meningkatkan motivasi belajar.
        \item Secara keseluruhan, rata-rata gabungan dari ahli materi, ahli media, uji coba kelas kecil, dan uji coba kelas besar adalah \(4{,}56\), termasuk kategori \textbf{Sangat Baik}. Hal ini menunjukkan bahwa media pembelajaran berbasis simulasi 3D layak digunakan untuk materi bangun ruang sisi datar kelas VIII.
    \end{enumerate}
\end{enumerate}
\textbf{Saran}\\
\addcontentsline{toc}{section}{B. Saran}
\hspace*{1cm}Berdasarkan hasil penelitian dan kesimpulan yang diperoleh, peneliti menyampaikan beberapa saran sebagai berikut:
\begin{enumerate}
    \item Media pembelajaran berbasis simulasi 3D ini diharapkan dapat menjadi alternatif bagi pendidik dalam menyampaikan materi bangun ruang sisi datar, serta memotivasi pengembangan media pembelajaran berbasis teknologi untuk materi lainnya.
    \item Penelitian selanjutnya dapat menguji efektivitas penggunaan media ini, pengaruhnya terhadap motivasi belajar peserta didik, serta dampaknya pada hasil belajar matematika.
    \item Teknologi 3D sebaiknya tidak hanya digunakan untuk penyajian materi, tetapi juga untuk evaluasi pembelajaran, misalnya melalui soal atau tugas berbasis simulasi 3D.
\end{enumerate}
